\chapter{Statistiques}
\begin{prerequis}
Calculer une fréquence et une moyenne simple; lire un tableau et un diagramme;
ordonner une liste de nombres.
\end{prerequis}
\begin{automatismes}
Pour $2,3,3,4,8$, calculer l'effectif, la moyenne, la médiane et l'étendue;
indiquer l'effet du remplacement de $8$ par $80$ sur ces indicateurs.
\end{automatismes}
\begin{activite}[Activité d'entrée -- Une seule valeur pour résumer?]
Comparer les séries $A=(8,8,8,8,8)$ et $B=(4,6,8,10,12)$. Elles ont la même
moyenne. Sont-elles pourtant semblables? Proposer une mesure de la dispersion,
puis confronter les propositions à la variance et à l'écart-type.
\end{activite}

%% ════════════════════════════════════════════════════════════════

\begin{anecdote}[La loi des grands nombres]
Jacob Bernoulli ($1654$--$1705$) d\'emontra la loi des grands nombres :
plus on r\'ep\`ete une exp\'erience al\'eatoire, plus la fr\'equence d'un
\'ev\'enement s'approche de sa probabilit\'e. Il mit $20$ ans \`a trouver la
preuve. Son neveu Nicolas dit : \og Cela n'\'etait d\'ej\`a pas inconnu aux
sots.\fg --- ce qui, peut-\^{e}tre, \'etait une exag\'eration.
\end{anecdote}

\begin{objectifs}
\begin{itemize}
  \item Organiser des donn\'ees statistiques~: variable discr\`ete et continue.
  \item Calculer et interpr\'eter les indicateurs de position~:
        moyenne, m\'ediane, mode.
  \item Calculer et interpr\'eter les indicateurs de dispersion~:
        \'etendue, variance, \'ecart-type.
  \item Construire et lire les repr\'esentations graphiques~:
        diagramme en b\^atons, diagramme circulaire, histogramme.
\end{itemize}
\end{objectifs}

\section{Statistiques descriptives}

\begin{defn}[Indicateurs de position et de dispersion]
Pour une s\'erie de valeurs $x_1, \ldots, x_n$ :
\begin{itemize}
  \item \textbf{Moyenne} : $\bar{x} = \dfrac{1}{n}\displaystyle\sum_{i=1}^n x_i$.
  \item \textbf{M\'ediane} $Q_2$ : valeur centrale une fois les donn\'ees tri\'ees.
  \item \textbf{Quartiles} $Q_1$, $Q_3$ : valeurs \`a $25\%$ et $75\%$.
  \item \textbf{\'Ecart-type} :
        $\sigma = \sqrt{\dfrac{1}{n}\displaystyle\sum_{i=1}^n (x_i - \bar{x})^2}$.
\end{itemize}
\end{defn}

\horsprogramme{La bo\^ite \`a moustaches et les quartiles de la figure suivante
ne font pas partie des indicateurs \`a ma\^itriser en T10.}
\begin{figure}[H]
\centering
\begin{tikzpicture}[xscale=0.55,yscale=1.2]
  \draw[BN,-Stealth] (0,0)--(21.5,0) node[right,font=\small] {valeurs};
  % Boîte à moustaches
  \draw[BN,very thick] (5,0.3)--(5,-0.3);  % min
  \draw[BN,very thick] (5,0)--(9,0);       % moustache gauche
  \draw[BN,very thick,fill=TQ!20] (9,-0.6) rectangle (15,0.6);  % boîte
  \draw[OR,very thick] (12,-0.6)--(12,0.6); % médiane
  \draw[BN,very thick] (15,0)--(20,0);     % moustache droite
  \draw[BN,very thick] (20,0.3)--(20,-0.3); % max
  % Labels
  \node[BN,font=\small] at (5,-1) {Min};
  \node[TQ,font=\small] at (9,-1) {$Q_1$};
  \node[OR,font=\small] at (12,-1) {$Q_2$};
  \node[TQ,font=\small] at (15,-1) {$Q_3$};
  \node[BN,font=\small] at (20,-1) {Max};
\end{tikzpicture}
\caption{Structure d'une bo\^ite \`a moustaches.}
\end{figure}

\horsprogramme{Approfondissement sur les quartiles, les fréquences cumulées et d'autres résumés d'une série.}
\section{Caract\'eristiques compl\'ementaires}

\begin{defn}[Mode et classe modale]
Le \textbf{mode} d'une s\'erie statistique est la valeur (ou la classe) qui
appara\^it le plus souvent. Dans une distribution en classes, on parle de
\textbf{classe modale}.
\end{defn}

\begin{defn}[Fr\'equences et fr\'equences cumul\'ees]
La \textbf{fr\'equence} d'une valeur $x_i$ est $f_i = n_i / N$.
La \textbf{fr\'equence cumul\'ee croissante} jusqu'\`a $x_k$ est
$F_k = \displaystyle\sum_{i \leq k} f_i$.
La m\'ediane est la valeur correspondant \`a une fr\'equence cumul\'ee de $50\%$.
\end{defn}


\section{S\'eries statistiques group\'ees}

\begin{defn}[Tableau de distribution]
\index{tableau de distribution}
Pour $n$ valeurs $x_1,\ldots,x_p$ aux effectifs $n_1,\ldots,n_p$
(avec $N=\sum n_i$), la \textbf{fr\'equence} de $x_i$ est $f_i=n_i/N$.
\end{defn}

\begin{defn}[Indicateurs de position]
\index{moyenne}\index{médiane}\index{mode}
\begin{itemize}
  \item \textbf{Moyenne~:} $\bar{x} = \dfrac{1}{N}\sum_{i=1}^p n_i x_i$.
  \item \textbf{M\'ediane~:} valeur $Q_2$ partageant la s\'erie en deux moiti\'es \'egales.
  \item \textbf{Mode~:} valeur (ou classe) de plus grande fr\'equence.
  \item \textbf{Quartiles~:} $Q_1$ (25\%), $Q_2$ (50\%), $Q_3$ (75\%).
\end{itemize}
\end{defn}

\begin{defn}[Indicateurs de dispersion]
\index{écart-type}\index{variance}
\begin{itemize}
  \item \textbf{\'Etendue~:} $e = x_{\max} - x_{\min}$.
  \item \textbf{Variance~:}
        $V = \dfrac{1}{N}\sum_{i=1}^p n_i(x_i-\bar{x})^2
           = \overline{x^2} - \bar{x}^2$
        (variance $=$ moyenne des carr\'es moins carr\'e de la moyenne).
  \item \textbf{\'Ecart-type~:} $\sigma = \sqrt{V}$.
\end{itemize}
\end{defn}

\begin{methode}[Calculer la m\'ediane sur donn\'ees group\'ees]
\index{médiane!calcul}
\begin{enumerate}
  \item Trier les valeurs dans l'ordre croissant.
  \item Calculer les effectifs cumul\'es.
  \item La m\'ediane est~:
    \begin{itemize}
      \item si $N$ est impair~: la valeur de rang $(N+1)/2$.
      \item si $N$ est pair~: la demi-somme des valeurs de rangs $N/2$ et $N/2+1$.
    \end{itemize}
\end{enumerate}
\end{methode}

\begin{figure}[H]
\centering
\begin{tikzpicture}[xscale=0.55,yscale=1.2]
  \draw[BN,-Stealth] (0,0)--(21.5,0) node[right,font=\small] {valeurs};
  \draw[BN,very thick] (5,0.3)--(5,-0.3);
  \draw[BN,very thick] (5,0)--(9,0);
  \draw[BN,very thick,fill=TQ!20] (9,-0.6) rectangle (15,0.6);
  \draw[OR,very thick] (12,-0.6)--(12,0.6);
  \draw[BN,very thick] (15,0)--(20,0);
  \draw[BN,very thick] (20,0.3)--(20,-0.3);
  \node[BN,font=\small] at (5,-1) {Min};
  \node[TQ,font=\small] at (9,-1) {$Q_1$};
  \node[OR,font=\small] at (12,-1) {$Q_2$};
  \node[TQ,font=\small] at (15,-1) {$Q_3$};
  \node[BN,font=\small] at (20,-1) {Max};
  \draw[TQ,<->,thick] (9,-1.6)--(15,-1.6)
    node[midway,below,font=\small] {Interquartile $Q_3-Q_1$};
\end{tikzpicture}
\caption{Bo\^ite \`a moustaches : cinq indicateurs.}
\end{figure}

\section{Exercices}

\subsection*{\diff{1}\quad Niveau 1 \textemdash\ Indicateurs de base}

\begin{exo}[Indicateurs sur une s\'erie \corrige]{1}
Notes d'une classe de 10 \'el\`eves~:
$8, 12, 14, 7, 15, 9, 13, 11, 16, 10$.
\begin{enumerate}
  \item Calculer la moyenne $\bar{x}$.
  \item Trier les donn\'ees et calculer la m\'ediane.
  \item Calculer $Q_1$ et $Q_3$.
  \item Calculer l'\'etendue et l'\'ecart-type $\sigma$.
\end{enumerate}
\end{exo}

\begin{exo}[Tableau de fr\'equences \corrige]{1}
R\'esultats d'un sondage sur les heures de t\'el\'evision par semaine~:
\begin{center}
\begin{tabular}{|c|ccccc|}
\hline
\rowcolor{BN!10}
Heures & 0 & 1 & 2 & 3 & 4 \\
\hline
Effectif & 3 & 8 & 10 & 6 & 3 \\
\hline
\end{tabular}
\end{center}
\begin{enumerate}
  \item Calculer le nombre total d'\'el\`eves.
  \item Calculer les fr\'equences (en \%).
  \item Calculer la moyenne.
  \item Quel est le mode~?
\end{enumerate}
\end{exo}

\begin{exo}[Histogramme en classes \corrige]{1}
Masses (kg) de colis~:
\begin{center}
\begin{tabular}{|c|c|c|c|c|}
\hline
\rowcolor{BN!10}
Classe & $[0;2[$ & $[2;4[$ & $[4;6[$ & $[6;8]$ \\
\hline
Effectif & 3 & 8 & 6 & 3 \\
\hline
\end{tabular}
\end{center}
\begin{enumerate}
  \item Construire l'histogramme.
  \item Calculer la moyenne (centre de classe).
  \item Quelle est la classe modale~?
  \item \'Estimer la m\'ediane.
\end{enumerate}
\end{exo}

\begin{exo}[Diagramme circulaire \corrige]{1}
Modes de transport vers le lyc\'ee~:
bus $120$, pied $80$, v\'elo $45$, voiture $35$.
\begin{enumerate}
  \item Calculer les fr\'equences en \%.
  \item Construire le diagramme circulaire (calculer les angles en degr\'es).
  \item Quel mode repr\'esente $43{,}3\%$~?
\end{enumerate}
\end{exo}

\begin{exo}[Compl\'eter un tableau \corrige]{1}
Pour une s\'erie de $20$ valeurs~:
\begin{center}
\begin{tabular}{|c|c|c|c|}
\hline
\rowcolor{BN!10}
Valeur $x_i$ & Effectif $n_i$ & Fr\'equence $f_i$ & $n_i\cdot x_i$ \\
\hline
$2$ & $4$ & & \\
$4$ & $7$ & & \\
$6$ & $5$ & & \\
$8$ & $?$ & & \\
\hline
Total & $20$ & $1$ & \\
\hline
\end{tabular}
\end{center}
\begin{enumerate}
  \item Compl\'eter le tableau.
  \item Calculer la moyenne.
  \item Calculer la m\'ediane.
\end{enumerate}
\end{exo}

\begin{exo}[Bo\^ite \`a moustaches \corrige]{1}
Pour les donn\'ees~: $2, 5, 7, 8, 10, 11, 14, 15, 18, 20$~:
\begin{enumerate}
  \item Calculer $Q_1$, $Q_2$, $Q_3$, Min, Max.
  \item Construire la bo\^ite \`a moustaches.
  \item Calculer l'\'ecart interquartile $Q_3-Q_1$.
\end{enumerate}
\end{exo}

\begin{exo}[Moyenne pond\'er\'ee \corrige]{1}
Notes d'un \'el\`eve~:
Math. $14$ (coef. 4), Physique $12$ (coef. 3), Fran\c cais $11$ (coef. 2),
Histoire $13$ (coef. 1).
\begin{enumerate}
  \item Calculer la moyenne pond\'er\'ee.
  \item Calculer la moyenne simple (non pond\'er\'ee).
  \item Comparer les deux moyennes~: laquelle favorise cet \'el\`eve~?
\end{enumerate}
\end{exo}

\begin{exo}[Lire un histogramme \corrige]{1}
Un histogramme repr\'esente la r\'epartition des \'ages dans un village~:
\begin{center}
\begin{tabular}{|c|c|c|c|c|c|}
\hline
\rowcolor{BN!10}
Classe d'\^age & $[0;20[$ & $[20;40[$ & $[40;60[$ & $[60;80[$ & $[80;100[$ \\
\hline
Fr\'equence & $18\%$ & $28\%$ & $30\%$ & $16\%$ & $8\%$ \\
\hline
\end{tabular}
\end{center}
\begin{enumerate}
  \item Le village compte $2500$ habitants. Calculer chaque effectif.
  \item Construire l'histogramme.
  \item Quelle est la classe modale~?
  \item Dans quelle classe se trouve la m\'ediane~? \'Estimer sa valeur.
\end{enumerate}
\end{exo}

\begin{exo}[Effectifs cumul\'es \corrige]{1}
\index{effectifs cumulés}
Pour la s\'erie~: $3, 5, 5, 7, 8, 8, 8, 10, 12, 14$~:
\begin{enumerate}
  \item Construire le tableau des effectifs et fr\'equences cumul\'ees.
  \item Tracer le polygone des fr\'equences cumul\'ees.
  \item Lire graphiquement $Q_1$, $Q_2$, $Q_3$.
\end{enumerate}
\end{exo}

\begin{exo}[S\'erie statistique~: analyse rapide \corrige]{1}
On compare deux s\'eries A et B~:
\begin{center}
\begin{tabular}{|c|cccccccccc|}
\hline
\rowcolor{BN!10} S\'erie A & 10 & 10 & 10 & 10 & 10 & 10 & 10 & 10 & 10 & 10 \\
S\'erie B & 5 & 6 & 7 & 9 & 10 & 10 & 11 & 13 & 14 & 15 \\
\hline
\end{tabular}
\end{center}
\begin{enumerate}
  \item Calculer la moyenne et l'\'ecart-type de chaque s\'erie.
  \item Laquelle est la plus r\'eguli\`ere~?
  \item L'\'ecart-type mesure-t-il bien la dispersion~?
\end{enumerate}
\end{exo}

\begin{exo}[Identifier un indicateur \corrige]{1}
Pour chaque situation, quel indicateur statistique est le plus pertinent~?
\begin{enumerate}
  \item On veut savoir si la \og valeur typique\fg\ d'un logement
        est repr\'esentative des prix du march\'e.
  \item On cherche le niveau \og moyen\fg\ d'une classe.
  \item On compare la r\'egularit\'e de deux machines de production.
  \item On cherche la taille \og la plus fr\'equente\fg\ d'une population.
\end{enumerate}
\end{exo}

\begin{exo}[Calculer la variance \corrige]{1}
Pour la s\'erie $5, 7, 8, 10, 10, 12$~:
\begin{enumerate}
  \item Calculer la moyenne $\bar{x}$.
  \item Calculer $\overline{x^2} = \frac{1}{n}\sum x_i^2$.
  \item D\'eduire la variance $V = \overline{x^2} - \bar{x}^2$.
  \item Calculer l'\'ecart-type $\sigma$.
\end{enumerate}
\end{exo}

\begin{exo}[R\`egle empirique \corrige]{1}
Pour une distribution \og bien r\'epartie\fg, environ $68\%$ des valeurs
sont dans $[\bar{x}-\sigma, \bar{x}+\sigma]$.
Pour la s\'erie des notes~:
$8, 9, 11, 12, 12, 13, 14, 15, 15, 16$~:
\begin{enumerate}
  \item Calculer $\bar{x}$ et $\sigma$.
  \item Calculer l'intervalle $[\bar{x}-\sigma, \bar{x}+\sigma]$.
  \item Compter combien de notes se trouvent dans cet intervalle.
  \item Calculer le pourcentage.
\end{enumerate}
\end{exo}

\begin{exo}[Diagramme en b\^atons \corrige]{1}
Les r\'esultats d'un lanceur de d\'e (20 essais)~:
\begin{center}
\begin{tabular}{|c|cccccc|}
\hline
\rowcolor{BN!10}
Face & 1 & 2 & 3 & 4 & 5 & 6 \\
\hline
Nb de fois & 2 & 4 & 5 & 3 & 3 & 3 \\
\hline
\end{tabular}
\end{center}
\begin{enumerate}
  \item Construire le diagramme en b\^atons.
  \item Calculer la moyenne des r\'esultats.
  \item Si le d\'e \'etait parfait, quelle serait la moyenne~?
\end{enumerate}
\end{exo}

\begin{exo}[Indicateurs~: r\'esum\'e \corrige]{1}
Pour chaque s\'erie, calculer TOUS les indicateurs (mode, m\'ediane,
moyenne, \'etendue, \'ecart-type) et construire la bo\^ite \`a moustaches~:
\begin{enumerate}
  \item $S_1 = \{1, 3, 3, 5, 7, 7, 7, 9\}$
  \item $S_2 = \{2, 4, 6, 8, 10, 12, 14, 16\}$
\end{enumerate}
Laquelle est la plus \og sym\'etrique\fg~? Laquelle est la plus dispers\'ee~?
\end{exo}

\subsection*{\diff{2}\quad Niveau 2 \textemdash\ Analyse et comparaisons}

\begin{exo}[Comparer deux s\'eries \corrige]{2}
Deux \'el\`eves comparent leurs notes sur 10 devoirs~:
\begin{center}
\begin{tabular}{|c|cccccccccc|}
\hline
\rowcolor{BN!10} Rova & 8 & 12 & 15 & 9 & 14 & 7 & 16 & 11 & 13 & 10 \\
\hline
Lova & 11 & 11 & 10 & 12 & 9 & 13 & 11 & 10 & 12 & 11 \\
\hline
\end{tabular}
\end{center}
\begin{enumerate}
  \item Calculer la moyenne, la m\'ediane et l'\'ecart-type de chaque s\'erie.
  \item Qui a les r\'esultats les plus r\'eguliers~? Justifier.
  \item Construire la bo\^ite \`a moustaches de chaque s\'erie.
\end{enumerate}
\end{exo}

\begin{exo}[Fr\'equences cumul\'ees et m\'ediane \corrige]{2}
Salaires mensuels (en milliers d'ariary)~:
\begin{center}
\begin{tabular}{|c|c|c|c|c|c|}
\hline
\rowcolor{BN!10}
Tranche & $[200;400[$ & $[400;600[$ & $[600;800[$ & $[800;1000[$ & $[1000;1200]$ \\
\hline
Effectif & 5 & 18 & 24 & 12 & 6 \\
\hline
\end{tabular}
\end{center}
\begin{enumerate}
  \item Calculer les fr\'equences et effectifs cumul\'es.
  \item Construire le polygone des effectifs cumul\'es.
  \item Estimer graphiquement la m\'ediane.
  \item Calculer la moyenne.
\end{enumerate}
\end{exo}

\begin{exo}[Analyse critique de statistiques \corrige]{2}
Un fabricant annonce~:
\og La moyenne de satisfaction de nos clients est $7{,}2/10$.\fg
\begin{enumerate}
  \item Si $200$ clients ont r\'epondu, la somme des notes est de combien~?
  \item Si l'\'ecart-type est $2{,}8$~: entre quelles valeurs se situent
        la plupart des r\'eponses~?
  \item Un deuxi\`eme fabricant a une moyenne de $6{,}8$ mais un \'ecart-type
        de $0{,}4$. Lequel pr\'ef\'erer~? Justifier.
  \item Pourquoi la m\'ediane serait-elle parfois plus pertinente
        que la moyenne~?
\end{enumerate}
\end{exo}

\begin{exo}[Variance~: propri\'et\'es \corrige]{2}
Soit une s\'erie $x_1,\ldots,x_n$ de moyenne $\bar{x}$ et d'\'ecart-type $\sigma$.
\begin{enumerate}
  \item D\'emontrer que $V = \overline{x^2} - \bar{x}^2$.
  \item Si on ajoute la constante $k$ \`a chaque valeur~: quelle est la
        nouvelle moyenne~? Le nouvel \'ecart-type~?
  \item Si on multiplie chaque valeur par $k$~: quelle est la nouvelle
        moyenne~? Le nouvel \'ecart-type~?
  \item Des notes sur $10$ ont $\bar{x}=6{,}5$ et $\sigma=1{,}8$.
        On les ramène sur $20$. Nouvelle moyenne et \'ecart-type.
\end{enumerate}
\end{exo}

\begin{exo}[S\'erie avec une inconnue \corrige]{2}
Une s\'erie de $6$ valeurs est $3, 5, 7, x, 9, 11$.
\begin{enumerate}
  \item Calculer la moyenne en fonction de $x$.
  \item Si la moyenne est $7$~: trouver $x$.
  \item Pour cette valeur de $x$~: calculer l'\'ecart-type.
  \item Si l'\'ecart-type est $\sqrt{6}$~: trouver les valeurs possibles de $x$.
\end{enumerate}
\end{exo}

\begin{exo}[Position de la moyenne par rapport \`a la m\'ediane \corrige]{2}
\begin{enumerate}
  \item Pour $S=\{1, 2, 3, 4, 100\}$~: calculer moyenne et m\'ediane.
        La moyenne repr\'esente-t-elle bien la s\'erie~?
  \item Pour $S=\{1, 1, 2, 2, 3, 100\}$~: m\^eme question.
  \item Dans quelles situations la m\'ediane est-elle plus repr\'esentative
        que la moyenne~?
  \item Construire une s\'erie de $7$ valeurs o\`u la m\'ediane est bien
        plus petite que la moyenne.
\end{enumerate}
\end{exo}

\begin{exo}[Bo\^ite \`a moustaches~: comparaison \corrige]{2}
Les bo\^ites \`a moustaches de deux groupes A et B ont les caract\'eristiques~:
\begin{center}
\begin{tabular}{|c|ccccc|}
\hline
\rowcolor{BN!10} & Min & $Q_1$ & $Q_2$ & $Q_3$ & Max \\
\hline
Groupe A & 2 & 5 & 8 & 12 & 20 \\
Groupe B & 4 & 7 & 9 & 11 & 15 \\
\hline
\end{tabular}
\end{center}
\begin{enumerate}
  \item Construire les deux bo\^ites.
  \item Quel groupe a la plus grande \'etendue~? Le plus grand \'ecart interquartile~?
  \item Le groupe B est-il plus concentr\'e~? Justifier.
  \item Quel groupe a le plus de valeurs sup\'erieures \`a $9$~?
\end{enumerate}
\end{exo}

\begin{exo}[Ind\'ependance de la moyenne et l'\'ecart-type \corrige]{2}
\begin{enumerate}
  \item Construire trois s\'eries de $5$ valeurs, chacune de moyenne $10$,
        mais d'\'ecart-types $0$, $2$ et $5$.
  \item Construire deux s\'eries de $5$ valeurs, l'une d'\'ecart-type $3$
        et de moyenne $5$, l'autre d'\'ecart-type $3$ et de moyenne $15$.
  \item La connaissance de la seule moyenne suffit-elle \`a d\'ecrire une s\'erie~?
        Et la seule \'ecart-type~?
\end{enumerate}
\end{exo}

\begin{exo}[S\'erie en classes~: \'estimer les quartiles \corrige]{2}
\index{quartiles!estimation}
Pour une s\'erie group\'ee en classes, on estime les quartiles par interpolation~:
si $Q_1$ se trouve dans la classe $[a;b[$ d'effectif $n_k$ avec un effectif
cumul\'e $C_{k-1}$ avant, alors~:
$Q_1 \approx a + \dfrac{N/4 - C_{k-1}}{n_k}\times(b-a)$.

\noindent
Appliquer \`a la distribution~:
\begin{center}
\begin{tabular}{|c|c|c|c|}
\hline
\rowcolor{BN!10}
Classe & $[10;15[$ & $[15;20[$ & $[20;25[$ \\
\hline
Effectif & 6 & 14 & 10 \\
\hline
\end{tabular}
\end{center}
\end{exo}

\begin{exo}[Graphique statistique~: choisir la repr\'esentation \corrige]{2}
Pour chaque situation, choisir et justifier la repr\'esentation graphique
la plus adapt\'ee (diagramme en b\^atons, circulaire, histogramme,
bo\^ite \`a moustaches)~:
\begin{enumerate}
  \item R\'epartition des \'el\`eves par option (LV1, LV2, option sport).
  \item Distribution des notes d'une classe (donn\'ees continues).
  \item Comparaison des r\'esultats de deux classes.
  \item \'Evolution du nombre de livres lus par mois sur un an.
\end{enumerate}
\end{exo}

\begin{exo}[Algorithme statistique \computer \corrige]{2}
\'Ecrire un programme Python qui~:
\begin{enumerate}
  \item Lit une liste de notes saisies par l'utilisateur.
  \item Calcule et affiche la moyenne, le minimum, le maximum,
        l'\'etendue et l'\'ecart-type.
  \item Calcule et affiche la m\'ediane (apr\`es tri).
  \item Indique si la moyenne est sup\'erieure ou inf\'erieure \`a $10$.
\end{enumerate}
\end{exo}

\begin{exo}[\logicoalgo\ Algorithme~: calculer m\'ediane et quartiles \corrige]{2}

\begin{algorithm}[H]
\caption{Calcul de la m\'ediane d'une s\'erie}
\KwIn{une liste $L$ de $n$ valeurs}
\KwOut{la m\'ediane}
Trier $L$ dans l'ordre croissant\;
\Si{$n \mod 2 = 1$}{
  \Retourner $L[(n+1)/2]$\tcp*[r]{rang du milieu (liste index\'ee depuis 1)}
}\Sinon{
  \Retourner $(L[n/2] + L[n/2+1]) / 2$\;
}
\end{algorithm}

\begin{enumerate}
  \item La condition \og$n$ est impair\fg\ est-elle \'equivalente \`a
        \og$n\mod2=1$\fg~? \'Enoncer la n\'egation de cette condition.
  \item Appliquer \`a $L=\{3,7,7,9,11\}$ (impair) et
        $L=\{2,5,7,8,10,12\}$ (pair).
  \item Modifier pour calculer aussi $Q_1$ (m\'ediane de la premi\`ere moiti\'e)
        et $Q_3$ (m\'ediane de la seconde moiti\'e).
  \item Impl\'ementer en Python \computer\ et v\'erifier sur les deux exemples.
\end{enumerate}
\end{exo}

\subsection*{\diff{3}\quad Niveau 3 \textemdash\ D\'emonstrations et analyse}

\begin{exo}[Variance et \'ecart-type~: propri\'et\'es formelles \corrige]{3}
Soit $f(t) = \dfrac{1}{n}\displaystyle\sum_{i=1}^n (x_i-t)^2$.
\begin{enumerate}
  \item D\'evelopper $f(t)$ et montrer que
        $f(t) = \overline{x^2} - 2t\bar{x} + t^2$.
  \item Montrer que $f$ est un trin\^ome en $t$ de coefficient positif.
  \item Montrer que $f$ est minimale pour $t=\bar{x}$.
  \item Conclure~: la variance est la valeur minimale de $f$.
\end{enumerate}
\end{exo}

\begin{exo}[Paradoxe de Simpson \corrige]{3}
\index{paradoxe de Simpson}
On compare les taux de r\'eussite de deux m\'ethodes A et B~:
\begin{center}
\begin{tabular}{|l|cc|cc|}
\hline
\rowcolor{BN!10} & \multicolumn{2}{c|}{Groupe 1} & \multicolumn{2}{c|}{Groupe 2} \\
\rowcolor{BN!10} & R\'eussis & Total & R\'eussis & Total \\
\hline
M\'ethode A & 9 & 10 & 60 & 100 \\
M\'ethode B & 80 & 100 & 5 & 10 \\
\hline
\end{tabular}
\end{center}
\begin{enumerate}
  \item Calculer le taux de r\'eussite de chaque m\'ethode dans chaque groupe.
  \item Quelle m\'ethode semble meilleure dans le groupe 1~? Dans le groupe 2~?
  \item Calculer le taux global de chaque m\'ethode (tous groupes confondus).
  \item Le r\'esultat est-il surprenant~? Expliquer le paradoxe.
\end{enumerate}
\end{exo}

\begin{exo}[Sensibilit\'e aux valeurs extr\^emes \corrige]{3}
\begin{enumerate}
  \item Pour $S = \{5, 6, 7, 8, 9\}$~: calculer $\bar{x}$ et $\sigma$.
  \item Remplacer la derni\`ere valeur par $90$~:
        recalculer $\bar{x}$ et $\sigma$.
  \item Remplacer \`a la place la m\'ediane par $90$~:
        calculer l'impact sur la m\'ediane et l'\'ecart-type.
  \item Conclure sur la robustesse relative de la m\'ediane et de la moyenne.
\end{enumerate}
\end{exo}

\begin{exo}[Corr\'elation \'el\'ementaire \corrige]{3}
On observe les paires $(x_i, y_i)$~:
$(1,2)$, $(2,4)$, $(3,5)$, $(4,4)$, $(5,6)$.
\begin{enumerate}
  \item Calculer $\bar{x}$, $\bar{y}$, $\sigma_x$, $\sigma_y$.
  \item Calculer $\overline{xy} = \frac{1}{n}\sum x_i y_i$.
  \item Le coefficient de corr\'elation lin\'eaire est
        $r = \dfrac{\overline{xy}-\bar{x}\bar{y}}{\sigma_x\sigma_y}$.
        Calculer $r$.
  \item Que signifie $r$ proche de $1$~? Proche de $0$~? Proche de $-1$~?
\end{enumerate}
\end{exo}

\begin{exo}[Construire une s\'erie \`a partir d'indicateurs \corrige]{3}
\begin{enumerate}
  \item Construire une s\'erie de $6$ valeurs de moyenne $8$,
        m\'ediane $8$, et \'ecart-type $\sqrt{6}$.
  \item Construire une s\'erie de $5$ valeurs de moyenne $10$,
        mode $10$, mais m\'ediane diff\'erente de $10$.
  \item Est-il possible d'avoir une s\'erie avec $\bar{x}=5$,
        $\sigma=0$, mais deux valeurs diff\'erentes~?
\end{enumerate}
\end{exo}

\begin{exo}[Statistiques bidimensionnelles~: r\'egression \corrige]{3}
\index{régression linéaire}
Pour la s\'erie bivariate $(x,y)$~: $(1,3)$, $(2,5)$, $(3,4)$, $(4,7)$, $(5,6)$~:
\begin{enumerate}
  \item Calculer $\bar{x}$, $\bar{y}$.
  \item La droite de r\'egression passe par $(\bar{x},\bar{y})$ et
        a pour pente $a=\dfrac{\sum(x_i-\bar{x})(y_i-\bar{y})}{\sum(x_i-\bar{x})^2}$.
        Calculer $a$.
  \item Trouver l'ordonnée \`a l'origine $b=\bar{y}-a\bar{x}$.
  \item \'Ecrire l'\'equation de la droite et pr\'edire $y$ pour $x=6$.
\end{enumerate}
\end{exo}

%─────────────────────────────────────────────────────────────────
\subsection*{Probl\`emes}
%─────────────────────────────────────────────────────────────────

\begin{probleme}[Analyse d'une classe \corrigeP]{1}
Les notes finales de $25$ \'el\`eves en math\'ematiques sont~:
$7, 8, 9, 9, 10, 10, 11, 11, 11, 12, 12, 12, 12, 13, 13, 14, 14, 15, 15, 15, 16, 17, 17, 18, 19$.
\begin{enumerate}
  \item Construire le tableau de distribution (effectif et fr\'equence).
  \item Calculer moyenne, m\'ediane, mode, \'ecart-type.
  \item Construire la bo\^ite \`a moustaches.
  \item Quel pourcentage d'\'el\`eves a obtenu une note sup\'erieure
        \`a la moyenne~? Inf\'erieure~?
\end{enumerate}
\end{probleme}

\begin{probleme}[Contr\^ole qualit\'e \corrigeP]{1}
Une machine fabrique des vis de longueur nominale $20$~mm.
On mesure $30$ vis~:
\begin{center}
\begin{tabular}{|c|c|c|c|c|c|}
\hline
\rowcolor{BN!10}
Longueur (mm) & 19,5 & 19,8 & 20,0 & 20,2 & 20,5 \\
\hline
Effectif & 3 & 7 & 12 & 6 & 2 \\
\hline
\end{tabular}
\end{center}
\begin{enumerate}
  \item Calculer la moyenne et l'\'ecart-type.
  \item Une vis est acceptable si $|l-20|\leq0{,}3$.
        Combien de vis sont rejet\'ees~?
  \item Calculer le pourcentage de rejet.
  \item Proposer une interpr\'etation de l'\'ecart-type en termes de qualit\'e.
\end{enumerate}
\end{probleme}

\begin{probleme}[\'Evolution d'une classe \corrigeP]{2}
En d\'ebut d'ann\'ee, une classe de $30$ \'el\`eves a une moyenne de $12$
et un \'ecart-type de $3$. En fin d'ann\'ee, un \'el\`eve de $18$ quitte
la classe et un \'el\`eve de $6$ arrive.
\begin{enumerate}
  \item Quelle est la somme des notes initiales~?
  \item Calculer la nouvelle moyenne.
  \item L'\'ecart-type peut-il augmenter ou diminuer~?
        Illustrer par un calcul.
\end{enumerate}
\end{probleme}

\begin{probleme}[Statistiques comparatives~: genre \corrigeP]{2}
Donn\'ees sur les salaires m\'edians en France (fictif, en milliers)~:
\begin{center}
\begin{tabular}{|c|cccccc|}
\hline
\rowcolor{BN!10}
Secteur & Commerce & Industrie & Sant\'e & IT & Education & BTP \\
\hline
Femmes & 22 & 28 & 30 & 40 & 26 & 24 \\
Hommes & 26 & 38 & 32 & 48 & 28 & 35 \\
\hline
\end{tabular}
\end{center}
\begin{enumerate}
  \item Calculer la moyenne des salaires f\'eminins et masculins.
  \item Calculer l'\'ecart-type de chaque distribution.
  \item Calculer le ratio F/H dans chaque secteur.
  \item Construire les deux bo\^ites \`a moustaches et les comparer.
\end{enumerate}
\end{probleme}

\begin{probleme}[Donn\'ees en classes~: \'estimation \corrigeP]{2}
La distribution des temps de trajet quotidiens (en minutes) d'un groupe~:
\begin{center}
\begin{tabular}{|c|c|c|c|c|c|}
\hline
\rowcolor{BN!10}
Classe & $[0;10[$ & $[10;20[$ & $[20;40[$ & $[40;60[$ & $[60;120]$ \\
\hline
Effectif & 8 & 15 & 20 & 10 & 7 \\
\hline
\end{tabular}
\end{center}
\begin{enumerate}
  \item Calculer la moyenne (centre de classe).
  \item Calculer l'\'ecart-type.
  \item Construire le polygone des fr\'equences cumul\'ees.
  \item Lire graphiquement $Q_1$, $Q_2$, $Q_3$.
\end{enumerate}
\end{probleme}

\begin{probleme}[Algorithme et statistiques \computer \corrigeP]{2}
\'Ecrire des programmes Python pour analyser une liste de donn\'ees~:
\begin{enumerate}
  \item Calculer la moyenne, le minimum, le maximum.
  \item Calculer la m\'ediane (apr\`es tri).
  \item Calculer l'\'ecart-type.
  \item Afficher un diagramme en b\^atons avec \texttt{matplotlib}
        (effectifs vs valeurs).
\end{enumerate}
Tester sur la s\'erie~: $4, 7, 7, 9, 11, 12, 14, 15, 15, 16$.
\end{probleme}

\begin{probleme}[Enquête r\'eelle~: concevoir et analyser \corrigeP]{1}
R\'ealiser une enquête dans la classe sur le temps consacr\'e aux r\'evisions
par semaine~:
\begin{enumerate}
  \item Proposer les modalit\'es de r\'eponse (classes d'heures).
  \item Collecter les donn\'ees de $20$ \'el\`eves (simuler si n\'ecessaire).
  \item Construire le tableau de distribution et l'histogramme.
  \item Calculer tous les indicateurs et commenter.
\end{enumerate}
\end{probleme}

\begin{probleme}[Taille de la bo\^ite et variation \corrigeP]{2}
Soit une s\'erie dont les $5$ indicateurs sont Min$=5$,
$Q_1=10$, $Q_2=15$, $Q_3=22$, Max$=30$.
\begin{enumerate}
  \item Construire la bo\^ite \`a moustaches.
  \item Si on ajoute $5$ \`a toutes les valeurs~:
        quels indicateurs changent, et de combien~?
  \item Si on multiplie toutes les valeurs par $2$~:
        quels indicateurs changent~?
  \item Calculer l'\'ecart interquartile dans chaque cas.
\end{enumerate}
\end{probleme}

\begin{probleme}[Comparaison internationale \corrigeP]{2}
L'esp\'erance de vie (ann\'ees) dans $8$ pays~: $68, 72, 75, 76, 79, 80, 81, 83$.
Pour les m\^emes pays, le revenu moyen (milliers de dollars)~:
$5, 8, 15, 18, 35, 42, 50, 60$.
\begin{enumerate}
  \item Calculer la moyenne et l'\'ecart-type de l'esp\'erance de vie.
  \item Calculer la moyenne et l'\'ecart-type du revenu.
  \item Tracer un diagramme de dispersion (points $(x_i, y_i)$).
  \item Semble-t-il y avoir une corr\'elation entre revenu et esp\'erance de vie~?
\end{enumerate}
\end{probleme}

\begin{probleme}[Suites et statistiques \corrigeP]{2}
On consid\`ere la suite $u_n = 2n+1$ pour $n = 1, 2, \ldots, 10$.
\begin{enumerate}
  \item Calculer les $10$ termes.
  \item Calculer la moyenne, la m\'ediane et l'\'ecart-type.
  \item Montrer que la moyenne est le terme du milieu.
  \item Calculer l'\'ecart-type en fonction du nombre de termes~:
        quel est le lien avec la raison de la suite~?
\end{enumerate}
\end{probleme}

\begin{probleme}[Transformation affine et statistiques \corrigeP]{2}
Si on applique la transformation $y_i = ax_i + b$ \`a chaque valeur~:
\begin{enumerate}
  \item Montrer que $\bar{y} = a\bar{x} + b$.
  \item Montrer que $\sigma_y = |a|\sigma_x$.
  \item Application~: des notes sur $10$ ont $\bar{x}=6$ et $\sigma=2$.
        On les ramène sur $20$. Calculer $\bar{y}$ et $\sigma_y$.
  \item Centrer-r\'eduire la s\'erie~: $z_i=(x_i-\bar{x})/\sigma$.
        Quelle est la moyenne de $z$~? Son \'ecart-type~?
\end{enumerate}
\end{probleme}

\begin{probleme}[Analyse de donn\'ees olympiques \corrigeP]{1}
Les temps (en secondes) des finalistes du $100$~m des J.O. 2020~:
$9{,}80$, $9{,}84$, $9{,}89$, $9{,}89$, $9{,}94$, $9{,}95$, $10{,}00$, $10{,}02$.
\begin{enumerate}
  \item Calculer la moyenne et l'\'ecart-type.
  \item Construire la bo\^ite \`a moustaches.
  \item Quelle est la diff\'erence entre le finaliste le plus rapide
        et le moins rapide~?
  \item Quelle est la diff\'erence entre la m\'ediane et la moyenne~?
\end{enumerate}
\end{probleme}

%─────────────────────────────────────────────────────────────────
\subsection*{D\'efis}
%─────────────────────────────────────────────────────────────────

\begin{challenge}[Paradoxe de Simpson]
\index{paradoxe de Simpson}
Dans un h\^opital, deux traitements $A$ et $B$~:
\begin{center}
\begin{tabular}{|l|cc|cc|}
\hline
\rowcolor{BN!10} & \multicolumn{2}{c|}{Femmes} & \multicolumn{2}{c|}{Hommes} \\
\rowcolor{BN!10} & Gu\'eris & Total & Gu\'eris & Total \\
\hline
$A$ & 200 & 250 & 45 & 50 \\
$B$ & 180 & 200 & 150 & 300 \\
\hline
\end{tabular}
\end{center}
\begin{enumerate}
  \item Calculer les taux de gu\'erison de $A$ et $B$ chez les femmes.
  \item M\^eme question chez les hommes.
  \item Calculer les taux globaux.
  \item Ce ph\'enom\`ene s'appelle le paradoxe de Simpson~:
        expliquer pourquoi les conclusions locales et globales divergent.
\end{enumerate}
\end{challenge}

\begin{challenge}[Variance minimale~: preuve]
\begin{enumerate}
  \item Montrer que parmi toutes les s\'eries de $n$ valeurs de m\^eme
        somme $S$, la s\'erie constante $x_i=S/n$ minimise la variance.
        \emph{(Utiliser la forme $V=\overline{x^2}-\bar{x}^2$ et
        l'in\'egalit\'e $\overline{x^2}\geq\bar{x}^2$.)}
  \item D\'emontrer l'in\'egalit\'e $\overline{x^2}\geq\bar{x}^2$.
        \emph{(C'est l'in\'egalit\'e de la variance positive.)}
  \item Quand a-t-on l'\'egalit\'e $V=0$~?
\end{enumerate}
\end{challenge}

\begin{challenge}[G\'en\'eration de distributions]
\begin{enumerate}
  \item Construire une s\'erie de $8$ valeurs avec $\bar{x}=10$,
        $\sigma=3$, m\'ediane$=10$, et une bo\^ite \`a moustaches sym\'etrique.
  \item Construire une s\'erie de $8$ valeurs avec $\bar{x}=10$
        mais une bo\^ite \`a moustaches fortement asym\'etrique.
  \item Peut-on construire une s\'erie avec $Q_1=Q_2=Q_3$~? Si oui, donner un exemple.
  \item Peut-on avoir $\sigma > Q_3-Q_1$~? Donner un exemple ou d\'emontrer l'impossibilit\'e.
\end{enumerate}
\end{challenge}

\begin{challenge}[Encadrement de la moyenne]
\begin{enumerate}
  \item Pour une s\'erie de $n$ valeurs dans $[m,M]$~: montrer que
        $m \leq \bar{x} \leq M$.
  \item Si $50\%$ des valeurs sont sup\'erieures \`a $Q_2$, est-ce que
        la moyenne est n\'ecessairement $\geq Q_2$~?
        Donner un exemple ou un contre-exemple.
  \item Pour une s\'erie sym\'etrique par rapport \`a sa moyenne
        (i.e.\ $\bar{x}=$ m\'ediane), est-ce que cela implique
        $Q_2-Q_1 = Q_3-Q_2$~?
\end{enumerate}
\end{challenge}

\begin{challenge}[Algorithme~: tri et calcul de quantiles \computer]
\begin{enumerate}
  \item \'Ecrire un programme Python qui trie une liste et calcule
        $Q_1$, $Q_2$, $Q_3$ automatiquement.
  \item Impl\'ementer la d\'etection de valeurs aberrantes~:
        une valeur est \og aberrante\fg\ si elle est \`a plus de $1{,}5\times(Q_3-Q_1)$
        de $Q_1$ ou $Q_3$.
  \item Tester sur une liste contenant une valeur aberrante.
  \item Comparer la moyenne avec et sans la valeur aberrante.
\end{enumerate}
\end{challenge}

\begin{challenge}[Statistiques bidimensionnelles : droite de r\'egression]
\index{droite de régression}
Cinq \'el\`eves ont obtenu les notes~: ($x$ = maths, $y$ = sciences)~:
$(8,10)$, $(12,13)$, $(15,14)$, $(10,11)$, $(18,17)$.
\begin{enumerate}
  \item Placer les points dans un rep\`ere.
  \item Calculer $\bar{x}$, $\bar{y}$, $\sigma_x$, $\sigma_y$.
  \item Calculer $\overline{xy}$ et la pente de la droite de r\'egression
        $a = (\overline{xy}-\bar{x}\bar{y})/\sigma_x^2$.
  \item \'Ecrire l'\'equation de la droite et l'utiliser pour pr\'edire
        la note en sciences d'un \'el\`eve qui a $13$ en maths.
\end{enumerate}
\end{challenge}

%─────────────────────────────────────────────────────────────────
\horsprogramme{La figure suivante associe un histogramme aux quartiles
d'une bo\^ite \`a moustaches~; seule la lecture de l'histogramme est exigible.}
\begin{figure}[H]
\centering
\begin{tikzpicture}[xscale=1.25,yscale=1.55]
  % Histogram
  \foreach \x/\h/\clr in {
    0/3/BN!40, 1/7/BN!55, 2/10/BN!70, 3/12/BN, 4/8/BN!70, 5/5/BN!55, 6/2/BN!40
  }{
    \fill[\clr] (\x,0) rectangle (\x+1,\h/14);
    \draw[BN!60,thin] (\x,0) rectangle (\x+1,\h/14);
  }
  % x-axis
  \draw[BN!50,->] (-0.3,0)--(7.5,0) node[right,font=\small] {valeurs};
  \draw[BN!50,->] (0,-0.05)--(0,1.05) node[above,font=\small] {};
  \foreach \x/\lbl in {0.5/10, 1.5/12, 2.5/14, 3.5/16, 4.5/18, 5.5/20, 6.5/22}{
    \node[below,font=\small,text=BN!70] at (\x,0) {\lbl};
  }
  \node[font=\large,text=BN] at (3.5,1.12) {Histogramme};
  % mean line
  \draw[OR,very thick,dashed] (3.5,0)--(3.5,0.88);
  \node[OR,font=\small,above] at (3.5,0.88) {$\bar{x}$};

  % Boxplot below
  \begin{scope}[yshift=-2.25cm]
    \draw[BN!40,->] (-0.3,0)--(7.5,0);
    % whiskers
    \draw[BN,very thick] (1,0)--(1.8,0);
    \draw[BN,very thick] (5.5,0)--(6.2,0);
    \draw[BN,very thick] (1,0.25)--(1,-0.25);
    \draw[BN,very thick] (6.2,0.25)--(6.2,-0.25);
    % box Q1..Q3
    \draw[BN,very thick,fill=TQ!15] (1.8,-0.35) rectangle (5.5,0.35);
    % median
    \draw[OR,very thick] (3.2,-0.35)--(3.2,0.35);
    % labels
    \node[below,font=\scriptsize,text=BN] at (1,-0.3) {minimum};
    \node[below,font=\scriptsize,text=TQ] at (1.8,-0.42) {$Q_1$};
    \node[below,font=\scriptsize,text=OR] at (3.2,-0.42) {$Q_2$};
    \node[below,font=\scriptsize,text=TQ] at (5.5,-0.42) {$Q_3$};
    \node[below,font=\scriptsize,text=BN] at (6.2,-0.3) {maximum};
    \node[font=\large,text=BN] at (3.5,0.78) {Bo\^ite \`a moustaches};
  \end{scope}
\end{tikzpicture}
\caption{Lire une distribution~: l'histogramme montre la forme g\'en\'erale~;
la bo\^ite \`a moustaches r\'esume cinq valeurs (approfondissement facultatif).
Ici la distribution est l\'eg\`erement asym\'etrique : la m\'ediane $Q_2$ (orange)
est \`a gauche de la moyenne $\bar{x}$ (orange, tiret\'e dans l'histogramme).}
\end{figure}

\begin{bilan}
\begin{itemize}
  \item \textbf{Indicateurs de position}~:
        m\'ediane (valeur centrale)~;
        moyenne $\bar{x}=\frac{1}{N}\sum n_i x_i$.
  \item \textbf{Indicateurs de dispersion}~:
        \'etendue $= x_{\max}-x_{\min}$~;
        variance $V=\overline{x^2}-\bar{x}^2$~;
        \'ecart-type $\sigma=\sqrt{V}$.
  \item \textbf{Repr\'esentations}~: diagramme en b\^atons (discret),
        circulaire (fr\'equences), histogramme (continu).
\end{itemize}

%─────────────────────────────────────────────────────────────────
\horsprogramme{Les exercices de logique et d'algorithmique statistique
ci-apr\`es peuvent mobiliser des quartiles, des simulations et des
transformations de s\'eries~: ils ne sont pas exigibles en T10.}
\subsection*{Logique \& Algorithmes}
%─────────────────────────────────────────────────────────────────

\begin{exo}[\logiq\ Propositions statistiques \corrige]{1}
\'Enoncer la n\'egation et d\'eterminer si vraie~:
\begin{enumerate}
  \item $\forall$ s\'erie, la moyenne est \'egale \`a la m\'ediane.
  \item $\exists$ s\'erie de $5$ valeurs de moyenne $10$ et d'\'ecart-type $0$.
  \item $\forall$ s\'erie, $Q_1\leq Q_2\leq Q_3$.
  \item L'\'ecart-type est toujours positif.
\end{enumerate}
\end{exo}

\begin{exo}[\logiq\ Implications entre indicateurs \corrige]{1}
Pour chaque paire, indiquer $\Rightarrow$, $\Leftarrow$ ou ni l'un ni l'autre~:
\begin{enumerate}
  \item \og $\sigma=0$\fg\ et \og toutes les valeurs sont \'egales\fg.
  \item \og La m\'ediane est plus grande que la moyenne\fg\ et
        \og la distribution est asym\'etrique\fg.
  \item \og L'\'etendue est nulle\fg\ et \og $\sigma=0$\fg.
  \item \og Le mode est unique\fg\ et \og la distribution est unimodale\fg.
\end{enumerate}
\end{exo}

\begin{exo}[\logiq\ Raisonnement par l'absurde en statistiques \corrige]{2}
D\'emontrer par l'absurde~:
\begin{enumerate}
  \item Une s\'erie ne peut pas avoir \`a la fois $\sigma>0$ et toutes ses valeurs \'egales.
  \item Si $\sigma=0$, alors $\bar{x}$ est la seule valeur de la s\'erie.
  \item Il est impossible que $Q_1>Q_3$.
\end{enumerate}
\end{exo}

\begin{exo}[\logiq\ Conditions sur la variance \corrige]{2}
\begin{enumerate}
  \item Montrer que $V=\overline{x^2}-\bar{x}^2\geq0$ pour toute s\'erie.
        Quel raisonnement~?
  \item La condition $\overline{x^2}=\bar{x}^2$ est-elle CN, CS ou CNS
        pour $\sigma=0$~?
  \item \'Enoncer~: \og La variance est nulle ssi toutes les valeurs sont \'egales.\fg
        D\'emontrer les deux sens.
\end{enumerate}
\end{exo}

\begin{exo}[\logiq\ Logique du paradoxe de Simpson \corrige]{3}
(Voir d\'efi Paradoxe de Simpson.)
\begin{enumerate}
  \item \'Enoncer le paradoxe~: \og Il peut exister des populations $G_1$, $G_2$
        et des traitements $A$, $B$ tels que $A$ soit meilleur que $B$ dans
        chaque groupe, mais $B$ soit meilleur que $A$ globalement.\fg
  \item Formuler cela comme une n\'egation d'une implication.
  \item D\'emontrer par un exemple num\'erique.
  \item Quelle en est la cause logique~? (pond\'erations diff\'erentes)
\end{enumerate}
\end{exo}

\begin{exo}[\logiq\ Propositions sur la bo\^ite \`a moustaches \corrige]{2}
\begin{enumerate}
  \item La proposition~: \og $50\%$ des donn\'ees sont dans la bo\^ite $[Q_1,Q_3]$\fg.
        Est-elle toujours vraie~? (Consid\'erer des donn\'ees \`a valeurs r\'ep\'et\'ees.)
  \item \og Si la bo\^ite est sym\'etrique, alors la moyenne \'egale la m\'ediane.\fg
        Vraie ou fausse~?
  \item Construire une s\'erie v\'erifiant $Q_1=Q_2$ mais $Q_2\ne Q_3$.
\end{enumerate}
\end{exo}

\begin{exo}[\algo\ Algorithme~: calculer moyenne et \'ecart-type \computer \corrige]{1}
(Voir exercice \'Etude de $f(x)=1/x$.)
\'Ecrire un programme Python qui calcule $\bar{x}$ et $\sigma$
\`a partir d'une liste en utilisant la formule $V=\overline{x^2}-\bar{x}^2$.
Tester sur $\{5,7,8,10,10,12\}$.
\end{exo}

\begin{exo}[\algo\ Algorithme~: tri par insertion \computer \corrige]{2}
\'Ecrire un algorithme de tri par insertion~:

\begin{algorithm}[H]
\caption{Tri par insertion}
\KwIn{liste $L$ de $n$ valeurs}
\Pour{$i \leftarrow 2$ \KwTo $n$}{
  $v \leftarrow L[i]$~;\quad $j \leftarrow i - 1$\;
  \TantQue{$j \geq 1$ \textbf{et} $L[j] > v$}{
    $L[j+1] \leftarrow L[j]$~;\quad $j \leftarrow j - 1$\;
  }
  $L[j+1] \leftarrow v$\;
}
\end{algorithm}

Tracer pour $L=\{5,3,8,1,4\}$. Impl\'ementer \computer.
\end{exo}

\begin{exo}[\algo\ Algorithme~: calcul des quartiles apr\`es tri \computer \corrige]{2}
(Voir exercice \logicoalgo\ dans ce chapitre.)
\'Etendre l'algorithme pour calculer aussi l'\'ecart-type, l'\'etendue,
et d\'etecter les valeurs aberrantes (plus de $1{,}5\times(Q_3-Q_1)$
hors de la bo\^ite).
Tester sur $\{2,5,7,8,9,10,11,14,50\}$.
\end{exo}

\begin{exo}[\algo\ Algorithme~: histogramme en classes \computer \corrige]{2}
\'Ecrire un programme Python qui, pour une liste de donn\'ees et
une amplitude de classe $h$~:
\begin{enumerate}
  \item Construit le tableau de distribution.
  \item Trace l'histogramme avec \texttt{matplotlib}.
  \item Calcule la fr\'equence cumul\'ee et la trace.
  \item Estime la m\'ediane graphiquement.
\end{enumerate}
\end{exo}

\begin{exo}[\algo\ Algorithme~: droite de r\'egression lin\'eaire \computer \corrige]{3}
(Voir d\'efi dans ce chapitre.)
\'Ecrire un programme Python complet qui~:
\begin{enumerate}
  \item Calcule $\bar{x}$, $\bar{y}$, $\sigma_x$, $\sigma_y$, et le coefficient $a$.
  \item Trace le nuage de points et la droite de r\'egression.
  \item Calcule le coefficient de corr\'elation $r$.
  \item Pr\'edit $y$ pour une valeur de $x$ hors de l'\'echantillon.
\end{enumerate}
\end{exo}

\begin{exo}[\algo\ Algorithme~: simulation d'une loi uniforme et moyenne \computer \corrige]{2}
\'Ecrire un programme Python qui simule $N$ tirages d'une variable al\'eatoire
uniforme sur $[0,10]$, calcule la moyenne et l'\'ecart-type, et v\'erifie
que la moyenne converge vers $5$ quand $N$ augmente.
Tester pour $N=10, 100, 1000, 10000$.
\end{exo}

%─────────────────────────────────────────────────────────────────
\subsection*{Probl\`emes Logique \& Algorithmes}
%─────────────────────────────────────────────────────────────────

\begin{probleme}[\logiq\ Logique du paradoxe de Simpson \corrigeP]{2}
(Voir d\'efi Paradoxe de Simpson.)
\begin{enumerate}
  \item Construire un exemple complet o\`u la m\'ethode $A$ est meilleure
        que $B$ dans chaque sous-groupe, mais pire globalement.
  \item \'Enoncer logiquement la condition qui garantit l'absence de paradoxe.
  \item La proposition~: \og Si $A$ est meilleure dans tous les sous-groupes,
        alors $A$ est meilleure globalement\fg\ est-elle toujours vraie~?
\end{enumerate}
\end{probleme}

\begin{probleme}[\logiq\ Logique de la transformation affine \corrigeP]{2}
Pour $y_i=ax_i+b$~:
\begin{enumerate}
  \item D\'emontrer que $\bar{y}=a\bar{x}+b$.
        Quelle propri\'et\'e logique de la somme utilise-t-on~?
  \item D\'emontrer que $\sigma_y=|a|\sigma_x$.
  \item La r\'eciproque~: si $\sigma_y=|a|\sigma_x$ et $\bar{y}=a\bar{x}+b$,
        peut-on conclure que $y_i=ax_i+b$~? Justifier.
\end{enumerate}
\end{probleme}

\begin{probleme}[\algo\ Algorithme~: analyse statistique compl\`ete \computer \corrigeP]{2}
\'Ecrire un programme Python complet qui, pour une s\'erie donn\'ee~:
\begin{enumerate}
  \item Calcule tous les indicateurs (mode, m\'ediane, moyenne, $\sigma$).
  \item D\'etecte les valeurs aberrantes.
  \item Trace l'histogramme et la bo\^ite \`a moustaches.
  \item Affiche un rapport r\'esum\'e.
\end{enumerate}
Tester sur les $25$ notes du probl\`eme d'analyse de classe.
\end{probleme}

\begin{probleme}[\algo\ Algorithme~: simulation et statistiques \computer \corrigeP]{2}
\'Ecrire un programme Python qui simule $1000$ lancers de deux d\'es
et calcule~:
\begin{enumerate}
  \item La distribution des sommes (de $2$ \`a $12$).
  \item La moyenne et l'\'ecart-type de cette distribution.
  \item L'histogramme des fr\'equences observ\'ees vs th\'eoriques.
  \item Le test~: la fr\'equence observ\'ee de chaque somme est-elle
        proche de la probabilit\'e th\'eorique~?
\end{enumerate}
\end{probleme}

%─────────────────────────────────────────────────────────────────
\subsection*{D\'efis Logique \& Algorithmes}
%─────────────────────────────────────────────────────────────────

\begin{challenge}[\logiq\ Variance minimale et in\'egalit\'e]
D\'emontrer que, pour toute s\'erie de $n$ valeurs~:
\begin{enumerate}
  \item $V\geq0$ (variance positive).
  \item $V=0 \Leftrightarrow$ toutes les valeurs sont \'egales.
  \item Parmi toutes les s\'eries de m\^eme somme $S=\sum x_i$, celle \`a variance
        minimale est la s\'erie constante.
  \item \'Enoncer tout cela sous forme d'implications et d'\'equivalences.
\end{enumerate}
\end{challenge}

\begin{challenge}[\algo\ Algorithme~: calculatrice statistique compl\`ete \computer]
\'Ecrire un programme Python interactif qui~:
\begin{enumerate}
  \item Permet \`a l'utilisateur de saisir des donn\'ees.
  \item Calcule et affiche tous les indicateurs.
  \item Propose plusieurs repr\'esentations graphiques.
  \item Permet de tester l'effet d'une transformation $y=ax+b$.
  \item G\'en\`ere un rapport exportable.
\end{enumerate}
\end{challenge}

\end{bilan}


%% ════════════════════════════════════════════════════════════════
\backmatter
